\documentclass{article}
\usepackage{amsmath,amssymb,geometry}
\geometry{a4paper,margin=1in}
\title{Finite Digest: Algebraic Metriplectic Bracket Packet}
\author{info-geometry-lean theorem-safe intake}
\date{}
\begin{document}
\maketitle

\section{Scope}
This digest records only a finite algebraic metriplectic interface.  It does not
construct a smooth manifold, a superKähler structure, a Fisher metric, or a
physical non-equilibrium thermodynamic model.

\section{Algebraic data}
Let $R$ be a commutative ring of observables.  A finite metriplectic packet
contains two brackets
\[
\{\cdot,\cdot\}:R\times R\to R,
\qquad
\langle\!\langle\cdot,\cdot\rangle\!\rangle:R\times R\to R,
\]
a Hamiltonian $H\in R$, and an entropy observable $S\in R$.  The Lean packet
stores the laws
\[
\{f,f\}=0,
\qquad
\langle\!\langle H,f\rangle\!\rangle=0,
\qquad
\langle\!\langle f,H\rangle\!\rangle=0,
\]
and
\[
\{S,f\}=0,
\qquad
\{f,S\}=0.
\]

\section{Evolution readouts}
The total algebraic evolution is
\[
\dot f=\{f,H\}+\langle\!\langle f,S\rangle\!\rangle.
\]
From the stored laws, Lean proves
\[
\dot H=0,
\qquad
\dot S=\langle\!\langle S,S\rangle\!\rangle.
\]
No positivity or second-law inequality is asserted without a concrete ordered
metric model.

\section{Artifacts}
\begin{itemize}
\item Lean: \texttt{lean/InfoGeometry/Topology/Metriplectic.lean}
\item GAP finite grading smoke test: \texttt{tools/gap/metriplectic\_superkahler.g}
\end{itemize}

\end{document}
